\documentclass[11pt,a4paper,spanish]{book}
\usepackage{estilo_unir-1}
\hypersetup{
	colorlinks   = true, %Colours links instead of ugly boxes
	urlcolor     = blue, %Colour for external hyperlinks
	linkcolor    = black, %Colour of internal links
	citecolor   = red %Colour of citations
}
\usepackage[nottoc]{tocbibind}

%---------------------------
%título de la asignatura, trabajo y autor
%---------------------------
\subject{Computación Cuántica}
\title{Escribir el título del documento}
\author{Enrique Prada Vázquez}
\profesor{Andrés Navas Cáliz}
\date{17-04-2026}

%---------------------------
%marges
%---------------------------
%\usepackage[margin=1.9cm]{geometry}
%---------------------------
%---------------------------
%---------------------------
%---------------------------
\begin{document}
\renewcommand{\listfigurename}{Índice de figuras}
\renewcommand{\listtablename}{Índice de tablas}
\renewcommand{\contentsname}{Índice de contenidos}
\renewcommand{\figurename}{Figura}
\renewcommand{\tablename}{Tabla} 

\maketitle
\frontmatter
\tableofcontents
\listoffigures
Fijaos que SOLO aparece el título de cada figura, no aparece toda la leyenda ni la fuente. \\
(QUITAR ESTA PÁGINA SI NO HAY FIGURAS) 
\listoftables
Fijaos que SOLO aparece el título de cada tabla, no aparece toda la leyenda ni la fuente. \\
(QUITAR ESTA PÁGINA SI NO HAY TABLAS)

\chapter{Resumen}
Este Trabajo presenta una simulación interactiva que visualiza estados cuánticos en la Esfera de Bloch. Utilizando la librería Qiskit para el motor de cálculo y Matplotlib para el renderizado gráfico, se ha construido una herramienta modular que representa la evolución de un cúbit ante diversas puertas lógicas. El Trabajo se centra en el mapeo de vectores de estado complejos a coordenadas cartesianas reales, permitiendo observar fenómenos de superposición y fase de manera intuitiva. Los resultados visualizan el comportamiento de operadores unitarios fundamentales, tales como la puerta de Hadamard y las puertas de Pauli.
\vspace{0.5cm}
{\bf Palabras clave:} computación cuántica, esfera de bloch, qiskit, cúbit, visualización de datos, mecánica cuántica.

\chapter{Abstract}
This Work presents an interactive simulation that visualizes quantum states on the Bloch Sphere. By utilizing the Qiskit library for the computational engine and Matplotlib for graphical rendering, a modular tool has been developed to represent the evolution of a qubit under various logic gates. This Work focuses on the mapping of complex state vectors to real Cartesian coordinates, allowing the observation of superposition and phase phenomena in an intuitive manner. The results visualize the behavior of fundamental unitary operators, such as the Hadamard gate and the Pauli gates.
\vspace{0.5cm}
{\bf Keywords:} quantum computing, bloch sphere, qiskit, qubit, data visualization, quantum mechanics.

\mainmatter

\chapter{Introducción}

El presente Trabajo aborda la simulación cúbits usando la representación geométrica de la Esfera de Bloch. La herramienta integra el cálculo de estados cuánticos mediante la librería Qiskit y una interfaz de visualización 3D programada en Matplotlib.

\section{Motivación y justificación}
El Trabajo se justifica como un ejercicio práctico para entender la manipulación de estados cuánticos y su representación geométrica, permitiendo asimilar conceptos como la superposición y la fase a través de la implementación directa del código.

\section{Planteamiento del Trabajo}
La implementación se basa en un diseño modular que separa la lógica de computación cuántica del renderizado gráfico. Se implementan funciones específicas para la preparación del estado, la aplicación de operadores unitarios mediante Qiskit y la proyección geométrica manual en un entorno tridimensional utilizando Matplotlib.

\section{Estructura del Trabajo}
La memoria se organiza en los siguientes apartados:
\begin{itemize}
    \item \textbf{Objetivos:} Detalle de las metas técnicas y específicas de la simulación.
    \item \textbf{Enunciado del problema:} Descripción de los requisitos y operaciones a implementar.
    \item \textbf{Desarrollo del Trabajo:} Explicación de la arquitectura del código, la conexión con Qiskit y el motor de visualización.
    \item \textbf{Resultados:} Ejecución y validación de la herramienta mediante casos de prueba estándar.
    \item \textbf{Conclusiones:} Resumen de los logros alcanzados y la funcionalidad del sistema.
\end{itemize}
\chapter{Objetivos}

\section{Objetivos generales}
El objetivo principal de este Trabajo es comprender las propiedades fundamentales del cúbit y su comportamiento dinámico mediante el desarrollo de una herramienta de simulación y visualización geométrica.

\section{Objetivos específicos}
\begin{itemize}
    \item Visualizar el impacto de las puertas cuánticas en la manipulación del vector de estado.
    \item Implementar un motor gráfico manual para proyectar estas propiedades físicas en la Esfera de Bloch.
    \item Validar el funcionamiento de los operadores unitarios mediante la observación de rotaciones en el espacio tridimensional.
\end{itemize}

\chapter{Enunciado del problema}

El problema central de este Trabajo consiste en la implementación de un sistema capaz de traducir el estado de un cúbit a una representación geométrica. Para ello, se deben resolver las siguientes cuestiones técnicas:

\begin{itemize}
    \item \textbf{Cálculo de la evolución cuántica:} Se requiere aplicar puertas lógicas sobre el vector de estado inicial, devolviendo el estado final resultante tras la operación.
    \item \textbf{Mapeo de coordenadas:} Transformar las amplitudes de probabilidad complejas $\alpha$ y $\beta$ del vector de estado en coordenadas reales $(x, y, z)$ dentro de un espacio euclídeo.
    \item \textbf{Renderizado manual:} Creación de un entorno gráfico tridimensional donde se visualice tanto la estructura de la Esfera de Bloch como el vector de estado y su orientación.
    \item \textbf{Representación de estados:} Con el fin de parametrizar el espacio de estados completo, el sistema debe permitir la entrada de parámetros $(\theta, \phi)$ para definir cualquier estado en la superficie de la esfera, gestionando correctamente la superposición y la fase.
\end{itemize}

En definitiva, se busca demostrar que mediante el uso de Qiskit y Matplotlib es posible construir un pipeline completo que ilustre la manipulación física de la información cuántica.


\chapter{Desarrollo del Trabajo}

El desarrollo se centra en la implementación de un sistema modular que integra el cálculo cuántico con una representación geométrica. A continuación, se detalla la arquitectura del simulador y la lógica de manipulación del cúbit.

\section{Arquitectura del Simulador}
El sistema se divide en dos bloques:
\begin{itemize}
    \item \textbf{Motor Cuántico (Qiskit):} Gestiona la creación de circuitos y la evolución del vector de estado mediante la aplicación de operadores.
    \item \textbf{Motor Gráfico (Matplotlib):} Renderiza manualmente la Esfera de Bloch y proyecta las amplitudes complejas en coordenadas cartesianas reales $(x, y, z)$.
\end{itemize}

\section{Manipulación mediante Puertas Cuánticas}
La parte central del desarrollo es la implementación de operadores unitarios que transforman el estado del cúbit. Se han integrado las siguientes familias de puertas:

\begin{itemize}
    \item \textbf{Puertas de Pauli (X, Y, Z):} Actúan como rotaciones de $\pi$ radianes sobre sus respectivos ejes. La puerta $X$ (NOT cuántico) realiza una inversión del bit, mientras que $Z$ modifica la fase relativa.
    \item \textbf{Puerta de Hadamard (H):} Es fundamental para la creación de superposición, transformando los estados base en estados de la base diagonal.
    \item \textbf{Puertas de Fase (S, T, Sdg):} Introducen rotaciones específicas en el eje $Z$ ($\pi/2$ y $\pi/4$ respectivamente), permitiendo ajustar la fase relativa del cúbit.
    \item \textbf{Puerta V (Raíz cuadrada de X):} Realiza una rotación de $\pi/2$ en el eje $X$, permitiendo situar el vector en posiciones intermedias entre los polos y el ecuador.
\end{itemize}

\section{Procesamiento y Proyección}
Tras la aplicación de las puertas, el sistema extrae el vector de estado resultante. Mediante la normalización de las amplitudes complejas, se calculan los componentes del vector de Bloch. Este proceso asegura que, independientemente de la secuencia de puertas aplicadas, el estado final sea representado con precisión geométrica sobre la superficie de la esfera, facilitando la interpretación visual de la computación cuántica.

\chapter{Conclusiones}

Tras el desarrollo e implementación del simulador de la Esfera de Bloch, se han alcanzado las siguientes conclusiones basadas en los objetivos planteados:

\begin{itemize}
    \item \textbf{Comprensión del estado cuántico:} El proceso de codificación manual del mapeo de amplitudes complejas a coordenadas reales ha permitido consolidar el entendimiento de la superposición y la fase. Se ha comprobado que la fase global no altera la posición del vector, mientras que la fase relativa determina su ubicación en el ecuador de la esfera.

    \item \textbf{Eficacia de la visualización geométrica:} La representación tridimensional mediante Matplotlib ha demostrado ser una herramienta indispensable para desmitificar el carácter abstracto de los operadores unitarios. Observar las rotaciones físicas ayuda a asimilar cómo las puertas cuánticas manipulan la información sin necesidad de recurrir únicamente al álgebra matricial.

    \item \textbf{Validación de operadores:} A través de los casos de prueba, se ha verificado que el motor de Qiskit responde con precisión a la teoría cuántica, permitiendo que puertas como Hadamard o la raíz cuadrada de X ($V$) lleven al cúbit a estados de superposición exacta.
\end{itemize}

En conclusión, este Trabajo ha servido para profundizar en la comprensión de los cúbits de manera intuitiva a través de un desarrollo que permite visualizar geométricamente sus distintos estados y cómo las puertas lógicas interactúan con ellos.


\begin{thebibliography}{99}
\bibitem[Lamport(1994)]{lamport1994} Lamport, L. (1994) \emph{\LaTeX: a document preparation system}, Addison
Wesley, Massachusetts, 2nd ed.
\bibitem[Ackerman(2017)]{ackerman2017} Ackerman, E. (2017) Why the English Premier League Should Have Playoffs.  McGraw Hill. London. 
\bibitem[Lund(2017)]{lund2017} Lund, A. P. and Bremner, Michael J. and Ralph, T. C. (2017) Quantum sampling problems, BosonSampling and quantum supremacy, npj Quantum Information 3: 15
\end{thebibliography}

\vspace{2cm}
Aún a riesgo de ser redundante, repito aquí que para consultar detenidamente la formas de referenciar, debéis seguir el formato APA del siguiente documento: \url{https://bibliografiaycitas.unir.net/documentos/APA_7ed_UNIR.pdf}

%\bibliographystyle{plain} 
%\bibliography{bibliografia}

\end{document}

